% generated by code/freeze_supp_labels.py from fqdoa_tc_supp.aux; do not edit
\newlabel{S-alg:fqdoa}{{S1}{1}{Search Procedure}{}{}}
\newlabel{S-sec:abs}{{S2}{1}{Absolute Cost of the Oracles}{}{}}
\newlabel{S-sec:other}{{S3}{1}{Other Decompositions}{}{}}
\newlabel{S-sec:thr}{{S4}{1}{Threshold Tests}{}{}}
\newlabel{S-tab:resources}{{S1}{2}{Cost of the whole search over $R$ iterations. Qubits: allocated (carrying gates). T-count: Proposition~\ref {M-prop:coherence}, confirmed by expanding every circuit. T-depth: standard seven-T Toffoli and four-T relative-phase Toffoli. Two-qubit: transpiled to $\{cx,id,rz,sx,x\}$ at optimization level 1 with all-to-all connectivity}{}{}}
\newlabel{S-tab:blocks}{{S2}{2}{T gates of one iteration by block, for D3, Lin, and the compile-time oracle (CT), from Proposition~\ref {M-prop:coherence}. The last row gives the two-qubit gates of the whole search, transpiled as in Table~\ref {tab:resources}, divided by $R$. n/a: the block does not exist}{}{}}
\newlabel{S-fig:scaling}{{S1}{2}{Cost of a full search against input width $n$ (Table~\ref {tab:resources}). (a) T-count, which dominates error-corrected cost: the coherent oracle for linear circuits (Lin) approaches CT. (b) Two-qubit gates after transpilation: the dense CNOT network of Lin grows faster, and Lin passes D3 at $n=14$. Logarithmic vertical axes}{}{}}
\newlabel{S-fig:sensitivity}{{S2}{2}{How the price of coherence varies with the function. (a) T-count ratio of D3 to CT for family $\famE $ at $n=\val {sens.n}$ against the number of exceptional cosets $c$, from Proposition~\ref {M-prop:coherence} with $k=\min (|\mathcal {M}|,c)$; dash-dotted: the lower bound of Proposition~\ref {M-prop:cheaper} for $|\mathcal {M}|=2$, which assumes $k=c$. (b) T-count ratio of Quad to CT per iteration against the number of Toffoli products $t$ in $U_f$. (c) Two-qubit gates per iteration of Lin, relative to D3 at the same $n$, against the density of $L$ (fraction of ones); dash-dotted: equal cost. Seeded random $L$; transpiled at optimization level 1 with all-to-all connectivity}{}{}}
\newlabel{S-tab:and}{{S3}{2}{T gates of one iteration when every multi-controlled gate is a ladder of logical AND gates, uncomputed by measurement. Twice: $U_f^{\dagger }$ evaluates $f$ a second time. Meas.: $U_f^{\dagger }$ is replaced by measuring $y$ and a phase correction (for D3, all $c$ pattern flips counted). Lin has no T gates in $U_f$, so measuring $y$ removes only its second CNOT network. Last two columns: ratio of D3 to CT}{}{}}
\newlabel{S-sec:classical}{{S5}{3}{Classical Control}{}{}}
\newlabel{S-sec:cons}{{S6}{3}{Construction Cost}{}{}}
\newlabel{S-sec:hwrec}{{S7}{3}{Hardware Records}{}{}}
\newlabel{S-sec:runmodel}{{\mbox  {S7-A}}{3}{Calibration Model of the Run}{}{}}
\newlabel{S-tab:thr}{{S4}{4}{Threshold tests $\mathcal {M}=\{y:y<\tau \}$ for families $\famN $ and $\famL $, benchmark instance. $k$: marked canonical inputs. T per iteration: the coherent oracle (Coh.), which marks with one MCZ per 1 bit of $\tau $, and the compile-time forms. These place one MCZ per marked input (Input), per cube (Cube), per term of the best fixed-polarity Reed--Muller form (RM), or per affine piece (Affine). Two-qubit: gates of the whole search, transpiled as in Table~\ref {tab:resources}, for the coherent oracle and for the compile-time form with the fewest T gates (Best CT). Choice: the pick of the selection pass for an error-corrected target, given every compile-time form; a tie goes to CT}{}{}}
\newlabel{S-sec:need}{{\mbox  {S7-B}}{4}{What Higher Success Would Require}{}{}}
\newlabel{S-sec:decisions}{{S9}{4}{Decisions of the Selection Pass}{}{}}
\newlabel{S-tab:thrR}{{S5}{5}{Threshold tests for family $\famR $, the quadratic layer of family $\famN $ followed by the linear layer of family $\famL $, benchmark instance. Columns as in Table~\ref {tab:thr}}{}{}}
\newlabel{S-tab:queries}{{S6}{5}{Oracle queries: classical random sampling of canonical inputs (expected evaluations of $f$) against Grover iterations $R$, with $k=2$}{}{}}
\newlabel{S-fig:hardware}{{S3}{5}{Success probability of every design (Table~\ref {M-tab:matched}), with the no-search level $k/2^{n-1}$ (dash-dotted) and the chance level $k/2^n$ (dotted) of each size. Triangles: FakeFez noise model with 95\% Wilson intervals. Diamonds: Eq.~\eqref {M-eq:survival} with the rate fitted to D3 at $n=4$ on \val {hw.A.date}. Squares and crosses: \texttt {ibm\_fez} and \texttt {ibm\_marrakesh}, means of \val {hw.runs_word} runs on \val {hw.B.date}. Subscripts give $R$ where it differs from the noiseless optimum; No: no-search circuit}{}{}}
\newlabel{S-sec:fake}{{S10}{5}{Validation on Device Models}{}{}}
\newlabel{S-tab:cons}{{S7}{6}{Measured construction cost on a laptop (\val {cons.cpus} logical processors, \val {cons.ram}~GB of memory), median of \val {cons.reps} runs. Listing: the marked inputs from the function description by the algorithms of Proposition~\ref {M-prop:listing}, slowest of the three families. Enumeration: every canonical input of a black box, per family, as the loop of HQ-DOA does. Build: the Qiskit circuit of the whole search. Transpile: to $\{cx,id,rz,sx,x\}$ at optimization level 1 with all-to-all connectivity. The process peaked at \val {cons.peak_mb}~MB}{}{}}
\newlabel{S-tab:hwruns}{{S8}{6}{Every hardware run: D3 alone on \texttt {ibm\_fez} on \val {hw.A.date}, and every design on \texttt {ibm\_fez} and \texttt {ibm\_marrakesh} on \val {hw.B.date}, when each job held all its circuits. Success: fraction of the \val {hw.shots} shots that return a marked canonical input}{}{}}
\newlabel{S-sec:period}{{S11}{6}{Cost of Finding the Period}{}{}}
\newlabel{S-sec:ablation}{{S12}{6}{Ablation of the Compilation Steps}{}{}}
\newlabel{S-tab:calib}{{S9}{7}{Calibration in force at the first run of each circuit, for the qubits and couplers it used. Columns: time of the calibration (UTC), and medians of the CZ, readout, and single-qubit ($sx$) errors, of the relaxation time $T_1$, and of the dephasing time $T_2$}{}{}}
\newlabel{S-tab:routing}{{S10}{7}{Two-qubit counts with and without routing to the Heron~r2 coupling map (FakeFez), both at optimization level 3 with native gates. Median over seeds, with the range in parentheses. Ratio: median of routed over unrouted. Run: counts of the circuits run on \texttt {ibm\_fez}; n/r: not run}{}{}}
\newlabel{S-tab:select}{{S11}{7}{Choice of the selection pass on the Heron~r2 model: design, iteration count, and predicted success. $^{\dagger }$Warning: no candidate is predicted to exceed the no-search level}{}{}}
\newlabel{S-tab:fake}{{S12}{8}{Validation of the selection pass on five device models: routed two-qubit count $G$, simulated success (Sim.), and predicted success (Pred.). D3 at $n=4$ calibrates the rate, so its prediction equals its simulation. The rows with $n=8$ extend the validation one size further, where only CT of the benchmark is simulated}{}{}}
\newlabel{S-tab:xover}{{S13}{8}{The family $\famR $ case of Section~\ref {M-sec:threshold} of the paper ($n=8$, $\tau =32$, $k=\val {xover.k}$) on five device models. $G$: routed two-qubit count (transpiler seed 0). Sim.: success under the device's noise model, with \val {xover.shots.ct} shots for CT and \val {xover.shots.coh} for the 22-qubit coherent oracle. Pred.: Eq.~\eqref {M-eq:survival} with the rate $p_{\mathrm {eff}}$ fitted on that device. CT is the affine cover, its cheapest synthesis. n/s: not simulated}{}{}}
\newlabel{S-tab:ablpass}{{S14}{8}{Ablation of the selection pass over the \val {seeds.decisions} decisions at $n=6$ (five device models, three transpiler seeds). Matches: decisions in which the rule picks the best simulated candidate. Regret: largest loss in simulated success}{}{}}
\newlabel{S-tab:seeds}{{S15}{8}{Decisions of the selection pass at $n=6$ for each device model and transpiler seed. $p_{\mathrm {eff}}$: rate fitted to D3 at $n=4$ routed with that seed. Choices: decisions in which the pass picks the best simulated candidate, out of six. Regret: largest loss in simulated success}{}{}}
\newlabel{S-tab:simon}{{S16}{9}{Cost of finding the period and of the whole pipeline, Eq.~\eqref {M-eq:pipeline}. $\Pr [y=0]$: exact probability of the uninformative Simon sample $y=0$. Mean samples for $\famE $: Monte Carlo over \val {simon.n4.trials} trials at $n=4$ down to \val {simon.n14.trials} at $n=14$, drawn from the exact distributions; for $\famL $ and $\famN $: exact expectation for samples uniform on the space orthogonal to $s$. Bound: Proposition~\ref {M-prop:simon}. Simon: $S\,T_S$, and Search: $R\,T_{\mathrm {it}}$ for D3 ($\famE $), Quad ($\famN $), Lin ($\famL $), and CT, in T gates}{}{}}
\newlabel{S-tab:ablation}{{S17}{9}{Ablation at $n=4$ with the coherent oracles for families $\famL $ (Lin) and $\famN $ (Quad) and the compile-time oracle (CT) for comparison. Qubits: allocated (carrying gates). Two-qubit counts after transpilation at optimization level 1. Success: exact density-matrix probability under per-gate depolarizing noise of strength $p$}{}{}}
